% id: 298-699
% book: 개념원리 대수 · 32 수학적 귀납법
% section: 연습문제 실력 UP
% figure: none
% answer: $3$
% answer_source: 답지
% variant_level: 0 (원본 전사)
% 이 파일은 build.mjs 가 items.json 에서 만든다. 고칠 때는 items.json 을 고친다.
\smash{\raisebox{4.5mm}{\dmkichul{개념원리 대수 298쪽 699번}}}
다음은 수열 $\{a_n\}$의 일반항 $a_n$이 $a_n=1+\dfrac{1}{2}+\dfrac{1}{3}+\cdots+\dfrac{1}{n}$일 때, $n\ge 2$인 모든 자연수 $n$에 대하여 등식 \par\smallskip\noindent\hspace*{1.5em}$n+a_1+a_2+a_3+\cdots+a_{n-1}=na_n$\hfill$\cdots\cdots$ ㉠\par\smallskip\noindent 이 성립함을 수학적 귀납법으로 증명하는 과정이다. \exprbox{\begin{minipage}{0.96\linewidth}\raggedright\lineskip=4pt {\small\bfseries 증명}\par\smallskip (i) $n=2$일 때, (좌변)$=3$, (우변)$=3$이므로 ㉠이 성립한다.\\ (ii) $n=k\ (k\ge 2)$일 때, ㉠이 성립한다고 가정하면\\ \hspace*{1.5em}$k+a_1+a_2+a_3+\cdots+a_{k-1}=ka_k$\\ 양변에 $\framebox[2.2em]{\scriptsize ㈎}$을 더하면\\ \hspace*{1.5em}$(k+1)+a_1+a_2+a_3+\cdots+a_{k-1}+a_k$\\ \hspace*{1.5em}$=ka_k+\framebox[2.2em]{\scriptsize ㈎}=(k+1)(a_{k+1}-\framebox[2.2em]{\scriptsize ㈏})+1$\\ \hspace*{1.5em}$=(k+1)a_{k+1}$\\ 따라서 $n=k+1$일 때에도 ㉠이 성립한다.\\ (i), (ii)에서 $n\ge 2$인 모든 자연수 $n$에 대하여 ㉠이 성립한다.\end{minipage}}\noindent 위의 과정에서 ㈎, ㈏에 알맞은 식을 각각 $f(k)$, $g(k)$라 할 때, $f(4)-g(11)$의 값을 구하시오.
